\documentclass[10pt]{article}
\usepackage[T1]{fontenc}
\usepackage[utf8]{inputenc}
\usepackage[french]{babel}
\usepackage{amsmath,amssymb,mathrsfs,mathpazo}
\usepackage{geometry}
\geometry{a4paper,margin=2cm}
\usepackage{graphicx,array,tabularx,tabularray,colortbl,multirow,multicol,booktabs,enumitem,pifont,ulem,fancybox,diagbox,cancel}
\usepackage{pstricks,pstricks-add,pst-plot,pst-tree,pst-node,pst-text,pst-func,pst-math,pst-bspline,pst-3dplot,pst-eucl,pst-solides3d}
\let\pscalculate\undefined
\usepackage{pst-bezier}
\usepackage{tikz,pgfplots,tkz-tab}
\usetikzlibrary{arrows,arrows.meta,patterns,calc,positioning,trees,decorations.markings,shapes.geometric,shapes.misc,shapes.arrows,intersections,decorations.pathmorphing,decorations.pathreplacing,backgrounds,shadows,babel}
\pgfplotsset{compat=1.18}
\usepgfplotslibrary{fillbetween}
\usepackage[np]{numprint}
\usepackage{listings,makecell,pas-tableur,twoopt}
\let\barre\undefined
\usepackage[pstricks]{tablvar}
\definecolor{darkgreen}{RGB}{0,100,0}
\providecommand{\red}{\color{red}}
\providecommand{\blue}{\color{blue}}
\providecommand{\green}{\color{green}}
\providecommand{\black}{\color{black}}
\providecommand{\white}{\color{white}}
\providecommand{\gray}{\color{gray}}
\providecommand{\cyan}{\color{cyan}}
\providecommand{\magenta}{\color{magenta}}
\providecommand{\yellow}{\color{yellow}}
\providecommand{\violet}{\color{violet}}
\providecommand{\orange}{\color{orange}}
\providecommand{\euro}{\texteuro}
\providecommand{\EUR}{\texteuro}
\providecommand{\ieme}{\textsuperscript{e}}
\providecommand{\degres}{\ensuremath{^\circ}}
\providecommand{\degree}{\ensuremath{^\circ}}
\providecommand{\vv}[1]{\overrightarrow{#1}}
\providecommand{\Coord}[1]{#1}
\providecommand{\decosix}{}
\providecommand{\starredbullet}{\ensuremath{\bullet}}
\pagestyle{empty}
\setlength{\parindent}{0pt}
\usepackage{adjustbox}\DeclareUnicodeCharacter{2236}{:}\DeclareUnicodeCharacter{2190}{\ensuremath{\leftarrow}}
\providecommand{\R}{\mathbb{R}}
\providecommand{\N}{\mathbb{N}}
\providecommand{\Z}{\mathbb{Z}}
\providecommand{\Q}{\mathbb{Q}}
\providecommand{\D}{\mathbb{D}}
\providecommand{\C}{\mathbb{C}}
\providecommand{\e}{{\mathrm{e}}}
\providecommand{\vect}[1]{\overrightarrow{#1}}
\providecommand{\vectt}[1]{\overrightarrow{\mathrm{#1}}}
\providecommand{\vvt}[1]{\overrightarrow{\mathrm{#1}}}
\providecommand{\np}[2][]{{#2~#1}}
\providecommand{\vv}[1]{{\overrightarrow{#1}}}
\renewcommand{\d}{{\mathrm{d}}}
\providecommand{\degres}{{^\circ}}
\providecommand{\degre}{{^\circ}}
\providecommand{\blue}{\color{blue}}
\providecommand{\red}{\color{red}}
\providecommand{\enskip}{\quad}
\providecommand{\llbracket}{\lbrack\!\lbrack}
\providecommand{\rrbracket}{\rbrack\!\rbrack}
\providecommand{\euro}{{\text{€}}}
\providecommand{\no}{{\text{n°}}}
\providecommand{\interval}[2]{{[#1;#2]}}
\providecommand{\intervFO}[2]{{[#1;#2[}}
\providecommand{\dotfill}{\ldots}

\providecommand{\ds}{\displaystyle}
\providecommand{\textding}[1]{#1}
\providecommand{\Oij}{\(\mathrm{O};\vec{\imath},\vec{\jmath}\)}
\providecommand{\Oijk}{\(\mathrm{O};\vec{\imath},\vec{\jmath},\vec{k}\)}
\providecommand{\Ouv}{\(\mathrm{O};\vec{u},\vec{v}\)}
\providecommand{\intervFF}[2]{\ensuremath{[#1;#2]}}
\renewcommand{\euro}{\ifmmode\text{\texteuro}\else\texteuro\fi}
\DeclareUnicodeCharacter{00B0}{\ifmmode^{\circ}\else\textdegree\fi}
\DeclareUnicodeCharacter{20AC}{\ifmmode\text{\texteuro}\else\texteuro\fi}
% Énoncés : Éducation nationale. Transcription et corrigés : APMEP, Denis Vergès et contributeurs. Découpage et classement : CoopMaths, Sébastien Lozano, Éric Elter et contributeurs. Adaptation de la présentation : Maths Explorer.
% CC BY-SA : https://coopmaths.fr/www/about/
% https://coopmaths.fr/alea/static/bac/2023/tex/bac_2023_08_sujet2_caledonie_1.tex
\begin{document}
\noindent\textit{APMEP / CoopMaths. CC BY-SA. Présentation adaptée par Maths Explorer.}\par\medskip
\medskip

\begin{minipage}{6.7cm}On considère le cube ABCDEFGH d'arête 1 représenté ci-contre.

On note K le milieu du segment [HG].

On se place dans le repère orthonormé $\left(\text{A}~;~\vect{\text{AD}},~\vect{\text{AB}}, \vect{\text{AE}}\right)$.

\smallskip

\begin{enumerate}
\item Justifier que les points C, F et K définissent un plan.
\item 
	\begin{enumerate}
		\item Donner, sans justifier, les longueurs KG, GF et GC.
		\item Calculer l'aire du triangle FGC.
		\item Calculer le volume du tétraèdre FGCK.
	\end{enumerate}
\end{enumerate}	
\end{minipage}\hfill
\begin{minipage}{6.7cm}
\psset{unit=0.75cm,arrowsize=2pt 3}
\begin{pspicture}(-2.8,-2.4)(5,5.3)
\psline[linestyle=dashed](3,0)(0,0)(0,3)
\psline[linestyle=dashed](-1.3,-1.3)
\psline[linewidth=1.25pt](3,0)(5,0)
\psline[linewidth=1.25pt](0,3)(0,5.3)
\psline[linewidth=1.25pt](-1.3,-1.3)(-2.8,-2.8)
\psframe[linewidth=1.25pt](-1.3,-1.3)(1.7,1.7)
\psline[linewidth=1.25pt](1.7,-1.3)(3,0)(3,3)(1.7,1.7)
\psline[linewidth=1.25pt](3,3)(0,3)(-1.3,1.7)
\uput[ur](0,0){A} \uput[dr](3,0){B} \uput[dr](1.7,-1.3){C} \uput[dr](-1.3,-1.3){D}
\uput[ul](0,3){E} \uput[ur](3,3){F} \uput[ul](1.7,1.7){G} \uput[ul](-1.3,1.7){H}
\uput[ul](-2.8,-2.8){$x$} \uput[u](5,0){$y$} \uput[r](0,5.3){$z$}
\psdots(0,0)(3,0)(1.7,-1.3)(-1.3,-1.3)(0,3)(3,3)(1.7,1.7)(-1.3,1.7)(0.2,1.7)
\uput[u](0.2,1.7){K}
\end{pspicture}
\end{minipage}

\smallskip
	
On rappelle que le volume $V$ d'un tétraèdre est donné par:
\[V = \dfrac13\mathcal{B} \times h,\]
où $\mathcal{B}$ est l'aire d'une base et $h$ la hauteur correspondante.

\begin{enumerate}[start=3]
\item 
	\begin{enumerate}
		\item On note $\vect{n}$ le vecteur de coordonnées $\begin{pmatrix}1\\2\\1\end{pmatrix}$.
		
Démontrer que $\vect{n}$ est normal au plan (CFK).
		\item  En déduire qu'une équation cartésienne du plan (CFK) est:
\[x + 2y + z - 3 = 0.\]
	\end{enumerate}
\item On note $\Delta$ la droite passant par le point G et orthogonale au plan (CFK). 

Démontrer qu'une représentation paramétrique de la droite $\Delta$ est:
\[\left\{\begin{array}{l c l}
x&=&1 + t\\
y&=&1 + 2t\\
z&=&1 + t
\end{array}\right. \quad  (t \in \R))\]
\item Soit L le point d'intersection entre la droite $\Delta$ et le plan (CFK).
	\begin{enumerate}
		\item Déterminer les coordonnées du point L.
		\item En déduire que LG $= \dfrac{\sqrt 6}{6}$.
	\end{enumerate}
\item En utilisant la question 2., déterminer la valeur exacte de l'aire du triangle CFK.
\end{enumerate}

\bigskip

\end{document}
